\section{APLICAÇÃO: DIAGNÓSTICO POR ANÚNCIO}

A Seção 5 mostra o que se associa à nota. Esta seção descreve como um anfitrião ou uma empresa gestora pode usar esse resultado. A ferramenta de diagnóstico gera, a pedido, o relatório de um único anúncio com quatro componentes.

O primeiro é o sentimento médio do anúncio em cada aspecto e a sua posição (percentil) entre os anúncios do mesmo bairro e do mesmo tipo de acomodação. Quando o bairro tem poucos anúncios, a comparação passa para o tipo de acomodação na cidade. O segundo componente é a evidência de cada aspecto. Aspectos com poucas sentenças recebem a marca de evidência baixa, sem rótulo de ponto forte ou de ponto de atenção, para que um número sustentado em duas frases não seja lido como diagnóstico. O terceiro componente reúne as queixas e os elogios recorrentes do anúncio, junto com a associação média, na cidade, desses tópicos com o percentil da nota (Seção 5). O quarto mostra a evolução do sentimento nas avaliações mais recentes. A ferramenta funciona também para anúncios novos: em vez de omitir um aspecto sem dados, marca a evidência como baixa.

Para a comparação foram construídas tabelas de referência com 66 grupos de bairro e tipo de acomodação que têm ao menos 30 anúncios e 3 grupos por tipo na cidade. Cada grupo traz a mediana e os quartis do sentimento por aspecto e da nota. O relatório respeita a classificação da Seção 4. Uma nota metodológica acompanha o documento e informa a precisão da atribuição por aspecto, e os aspectos que pedem cautela não devem orientar decisões isoladamente. Os aspectos classificados como confiáveis são limpeza, localização e custo-benefício.

Foram gerados três exemplos anonimizados: um anúncio de nota alta, um de nota baixa e um com poucas avaliações. No de nota baixa, o relatório aponta comodidades e satisfação geral muito abaixo dos pares e destaca as queixas de banheiro e cozinha. No de poucas avaliações, a maior parte dos aspectos aparece com evidência baixa. O relatório indica onde investigar. Ele não prevê a nota do anúncio e não prescreve causas, porque as associações da Seção 5 são médias da cidade e não efeitos sobre um anúncio específico.
